%% kapa: exam=1 round=36 year=2025 period=1 form=A subject=economics subject_ko=경제학원론 seq=33 number=73
%% source: Q-Net 감정평가사 기출문제 (cst00309) artlSeq=5231525 file=raw/1008/5231525/2025년 제36회 감정평가사 1차시험 1교시 문제지 원본.pdf
%% conversion: kordoc 4.24.51 -> tools/gen1.py (auto); review=proofread-sonnet; flags=none
\begin{question}{73}
국내총생산의 측정에 관한 설명으로 옳지 않은 것은?
\begin{choices}
  \item 시장에서 거래되지 않는 재화와 서비스는 포함되지 않는다.
  \item 국민총생산에서 국외순수취요소소득을 뺀 것과 같다.
  \item 일정기간 내에 생산됐지만 아직 판매되지 못한 최종재도 포함된다.
  \item 일정기간 국내에서 새로 생산된 재화와 서비스의 부가가치의 합으로 계산한다.
  \item 간접세와 감가상각도 포함된다.
\end{choices}
\end{question}
\answer{1}
